% from ega4-11-fr.tex
\oldpage[IV-3]{121}

\begin{lemma}[11.2.2.2]
\label{IV.11.2.2.2-fr}
Soient $(A_\alpha,g_{\beta\alpha})$ un système inductif filtrant d'anneaux,
$(M_\alpha,h_{\beta\alpha})$ un système inductif d'ensembles,
$A'=\varinjlim A_\alpha$, $M'=\varinjlim M_\alpha$,
$g_\alpha:A_\alpha\to A'$, $h_\alpha:M_\alpha\to M'$ les applications
canoniques. Pour tout $\alpha$, soit $F(M_\alpha)$ le $A_\alpha$-module des
combinaisons linéaires formelles d'éléments de $M_\alpha$~; soit de même
$F(M')$ le $A'$-module des combinaisons linéaires formelles d'éléments de
$M'$~; si $h^0_{\beta\alpha}:F(M_\alpha)\to F(M_\beta)$ (pour
$\alpha\leq\beta$) et $h'_\alpha:F(M_\alpha)\to F(M')$ sont les
$A_\alpha$-homomorphismes déduits de $h_{\beta\alpha}$ et $h_\alpha$
respectivement, $(F(M_\alpha),h^0_{\beta\alpha})$ est un système inductif
de $A_\alpha$-modules et $(h'_\alpha)$ un système inductif
d'homomorphismes. Alors
\[
h''=\varinjlim h'_\alpha:\varinjlim F(M_\alpha)\to
F(M')=F(\varinjlim M_\alpha)
\]
est un isomorphisme.
\end{lemma}

Pour la démonstration, voir Bourbaki, \emph{Alg.}, chap.~II, 3$^e$ éd.,
\S~6, n$^\circ$~6, cor. de la prop.~10.

\begin{env}[11.2.3]
\label{IV.11.2.3-fr}
Reprenons les notations de (11.2.1) et considérons particulièrement le cas
$i=1$~; posons $T=\operatorname{Tor}_1^A(M,N)$, $T'=T\otimes_AA'$,
$T''=\operatorname{Tor}_1^{A'}(M',N')$~; alors $\psi_1$ est
l'homomorphisme
\[
\operatorname{Ker}(f_0\otimes1_N)/\operatorname{Im}(f_1\otimes1_N)
\to
\operatorname{Ker}(f''_0\otimes1_{N'})/
\operatorname{Im}(f''_1\otimes1_{N'})
\]
qui se déduit par passage aux quotients de la restriction
$\operatorname{Ker}(f_0\otimes1_N)\to
\operatorname{Ker}(f''_0\otimes1_{N'})$ de
\[
u_1\otimes1:L_0\otimes_AN\to L''_0\otimes_{A'}N'=L''_0\otimes_AN.
\]
Posons $R=\operatorname{Ker}(\varepsilon)$,
$R''=\operatorname{Ker}(\varepsilon'')$, de sorte que l'on a les suites
exactes
\[
0\to R\xrightarrow{j}L_0\xrightarrow{\varepsilon}M\to0
\qquad\text{et}\qquad
0\to R''\xrightarrow{j''}L''_0\xrightarrow{\varepsilon''}M'\to0,
\]
d'où on déduit les suites exactes d'homologie
\begin{equation}
0=\operatorname{Tor}_1^A(L_0,N)\to T\xrightarrow{\partial}R\otimes_AN
\xrightarrow{j\otimes1}L_0\otimes_AN\xrightarrow{\varepsilon\otimes1}
M\otimes_AN\to0
\label{IV.11.2.3.1-fr}
\tag{11.2.3.1}
\end{equation}
\begin{equation}
0=\operatorname{Tor}_1^{A'}(L''_0,N')\to T''\xrightarrow{\partial''}
R''\otimes_{A'}N'\xrightarrow{j''\otimes1}L''_0\otimes_{A'}N'
\xrightarrow{\varepsilon''\otimes1}M'\otimes_{A'}N'\to0.
\label{IV.11.2.3.2-fr}
\tag{11.2.3.2}
\end{equation}
On a d'autre part un homomorphisme de $A'$-modules
\[
v:R'=\operatorname{Ker}(\varepsilon)\otimes_AA'
\to\operatorname{Ker}(\varepsilon')\xrightarrow{u_0}
\operatorname{Ker}(\varepsilon'')=R''.
\]
\end{env}

